
This section presents the planned data analytics pipeline for ingesting, processing, and
serving the Lichess dataset to two classes of consumer: (1) Analysts performing ad hoc
queries and generating reports via SQL/notebooks and (2) An
automated machine learning pipeline for supervised model training and serving.


\begin{figure}[htbp]
  \centering
  \includegraphics[width=\linewidth,keepaspectratio]{whowhatandhow.pdf}
  \caption{Pipeline stages: What, Why, Who, and How for each stage of the MLOps workflow.}
  \label{fig:pipeline-whowhatandhow}
\end{figure}


\begin{figure}[htbp]
  \centering
  \includegraphics[width=\linewidth,height=0.92\textheight,keepaspectratio]{arch_diag.pdf}
  \caption{End-to-end MLOps pipeline divided in terms of configuration, artefacts, and components.}
  \label{fig:mlops-pipeline}
\end{figure}

\section{Data ingestion}
\label{sec:stage-ingestion}

\begin{figure}[htbp]
  \centering
  \includegraphics[width=\linewidth,height=0.85\textheight,keepaspectratio]{dataIngestion.pdf}
  \caption{Data ingestion (\emph{extract} and \emph{load}): the orchestrator applies ingestion config, downloads the monthly \texttt{.pgn.zst} from Lichess to a raw-data bucket in MinIO, and validates checksums before landing the blob unchanged; checksum failures trigger alerts.}
  \label{fig:data-ingestion}
\end{figure}

This pipeline automates the ingestion of Lichess's monthly public game data using the Apache Airflow orchestrator. This stage covers the \emph{extract} and \emph{load} steps of the ELT pattern (\cref{fig:elt-flow}; \cref{ch:storage-planning}): the monthly shard is downloaded over HTTPS, its integrity is verified (SHA256), and the compressed file is landed in MinIO as a single immutable object with no parsing or reshaping.

Decompression, PGN parsing, and curated star-schema Parquet are performed in preprocessing (\cref{sec:stage-preprocessing}). We use MinIO as an S3-compatible store for raw and processed data; Parquet partitions in the processed bucket feed MariaDB and downstream analysts and ML tooling.

\section{Data pre-processing}
\label{sec:stage-preprocessing}

In preprocessing, a PySpark job implements the \emph{transform} step: it reads the raw \texttt{.zst} stream from MinIO, decompresses the Zstandard payload, parses each game into header-derived fields (and optionally move rows), builds the star-schema DataFrames (\texttt{fact\_games}, dimensions, optionally \texttt{fact\_moves}), and writes partitioned Parquet to the processed bucket.

After Spark writes those star‑schema Parquet files to MinIO, they are loaded into MariaDB, the columnar database which offers fast read. The tables containing the dimensions and the fact games are stored in MariaDB, so that fast SQL can be used for both validation. Great Expectations checkpoint is performed on the whole MariaDB table (or the corresponding Parquet) prior to any split. Verifies columns, data types, null ratios, and ranges of data. The pipeline stops if the validation fails and an alert is sent otherwise, an artefact is created which certifies the data quality. Feast (with MariaDB as its offline store) supports the point in time feature retrieval for the entire historical data. The time-dependent features (e.g., win rate of the players in the last 30 days) are calculated based on only data prior to the game's timestamp to prevent temporal leakage. This will create a feature DataFrame, which now includes all features and the target, split chronologically (e.g., 80 \% train, 20 \% test). Since features were created on the entire data set before the split, both partitions have all columns and there is no problem with “columns missing in training/test”.

\begin{figure}[htbp]
  \centering
  \includegraphics[width=\linewidth,height=0.92\textheight,keepaspectratio]{dataValidation.pdf}
  \caption{Data preprocessing under orchestration: train/test splits are checked for schema completeness and drift; Great Expectations stores validation reports; failures route to alerts and updated artefact keys, while validated processed data feeds an offline feature store.}
  \label{fig:data-validation}
\end{figure}


\section{Model development}
\label{sec:stage-model-dev}

\begin{figure}[htbp]
  \centering
  \includegraphics[width=\linewidth,height=0.85\textheight,keepaspectratio]{modelTrainer.pdf}
  \caption{Model training: load train/test arrays from validated artefacts, fit the model, compute metrics, register the best run in MLflow, and retain a containerised, reproducible training environment.}
  \label{fig:plan-model-trainer}
\end{figure}
The model training stage reads  data from database. Here, It uses the Feast for defining the feature from star schema. 

\section{Model deployment}
\label{sec:stage-deployment}

\begin{figure}[htbp]
  \centering
  \includegraphics[width=\linewidth,height=0.85\textheight,keepaspectratio]{deployment.pdf}
  \caption{Deployment: promote a recent model from MLflow, build a Dockerised FastAPI image, run health checks, and push to the serving host (e.g.\ EC2); improved metrics can trigger CI/CD to refresh the endpoint.}
  \label{fig:plan-deployment}
\end{figure}

Containers promote only after gates; MLOps owns CI/CD with SecOps image review (\cref{fig:plan-deployment}).
Canary and rollback retag registries; contracts keep serving aligned with preprocessing.

\section{Model monitoring}
\label{sec:stage-monitoring}

\begin{figure}[htbp]
  \centering
  \includegraphics[width=\linewidth,height=0.85\textheight,keepaspectratio]{monitoringObserv.pdf}
  \caption{Monitoring and observability: API metrics scraped by Prometheus, dashboards and automated reports, rule-based routing to Slack/PagerDuty/email, and controlled reruns from user or scheduled inputs.}
  \label{fig:plan-monitoring}
\end{figure}

Production tracks feature drift, latency SLOs, and score error as new months arrive (\cref{fig:plan-monitoring}).
MLOps triages alerts; retrains stay lineage-linked to ingestion batches.
